\IfFileExists{stacks-project.cls}{%
\documentclass{stacks-project}
}{%
\documentclass{amsart}
}

% For dealing with references we use the comment environment
\usepackage{verbatim}
\newenvironment{reference}{\comment}{\endcomment}
%\newenvironment{reference}{}{}
\newenvironment{slogan}{\comment}{\endcomment}
\newenvironment{history}{\comment}{\endcomment}

% For commutative diagrams we use Xy-pic
\usepackage[all]{xy}

% We use 2cell for 2-commutative diagrams.
\xyoption{2cell}
\UseAllTwocells

% We use multicol for the list of chapters between chapters
\usepackage{multicol}

% This is generally recommended for better output
\usepackage{lmodern}
\usepackage[T1]{fontenc}

% For cross-file-references
\usepackage{xr-hyper}

% Package for hypertext links:
\usepackage{hyperref}

% For any local file, say "hello.tex" you want to link to please
% use \externaldocument[hello-]{hello}
\externaldocument[introduction-]{introduction}
\externaldocument[conventions-]{conventions}
\externaldocument[sets-]{sets}
\externaldocument[categories-]{categories}
\externaldocument[topology-]{topology}
\externaldocument[sheaves-]{sheaves}
\externaldocument[sites-]{sites}
\externaldocument[stacks-]{stacks}
\externaldocument[fields-]{fields}
\externaldocument[algebra-]{algebra}
\externaldocument[brauer-]{brauer}
\externaldocument[homology-]{homology}
\externaldocument[derived-]{derived}
\externaldocument[simplicial-]{simplicial}
\externaldocument[more-algebra-]{more-algebra}
\externaldocument[smoothing-]{smoothing}
\externaldocument[modules-]{modules}
\externaldocument[sites-modules-]{sites-modules}
\externaldocument[injectives-]{injectives}
\externaldocument[cohomology-]{cohomology}
\externaldocument[sites-cohomology-]{sites-cohomology}
\externaldocument[dga-]{dga}
\externaldocument[dpa-]{dpa}
\externaldocument[sdga-]{sdga}
\externaldocument[hypercovering-]{hypercovering}
\externaldocument[schemes-]{schemes}
\externaldocument[constructions-]{constructions}
\externaldocument[properties-]{properties}
\externaldocument[morphisms-]{morphisms}
\externaldocument[coherent-]{coherent}
\externaldocument[divisors-]{divisors}
\externaldocument[limits-]{limits}
\externaldocument[varieties-]{varieties}
\externaldocument[topologies-]{topologies}
\externaldocument[descent-]{descent}
\externaldocument[perfect-]{perfect}
\externaldocument[more-morphisms-]{more-morphisms}
\externaldocument[flat-]{flat}
\externaldocument[groupoids-]{groupoids}
\externaldocument[more-groupoids-]{more-groupoids}
\externaldocument[etale-]{etale}
\externaldocument[chow-]{chow}
\externaldocument[intersection-]{intersection}
\externaldocument[pic-]{pic}
\externaldocument[weil-]{weil}
\externaldocument[adequate-]{adequate}
\externaldocument[dualizing-]{dualizing}
\externaldocument[duality-]{duality}
\externaldocument[discriminant-]{discriminant}
\externaldocument[derham-]{derham}
\externaldocument[local-cohomology-]{local-cohomology}
\externaldocument[algebraization-]{algebraization}
\externaldocument[curves-]{curves}
\externaldocument[resolve-]{resolve}
\externaldocument[models-]{models}
\externaldocument[functors-]{functors}
\externaldocument[equiv-]{equiv}
\externaldocument[pione-]{pione}
\externaldocument[etale-cohomology-]{etale-cohomology}
\externaldocument[proetale-]{proetale}
\externaldocument[relative-cycles-]{relative-cycles}
\externaldocument[more-etale-]{more-etale}
\externaldocument[trace-]{trace}
\externaldocument[crystalline-]{crystalline}
\externaldocument[spaces-]{spaces}
\externaldocument[spaces-properties-]{spaces-properties}
\externaldocument[spaces-morphisms-]{spaces-morphisms}
\externaldocument[decent-spaces-]{decent-spaces}
\externaldocument[spaces-cohomology-]{spaces-cohomology}
\externaldocument[spaces-limits-]{spaces-limits}
\externaldocument[spaces-divisors-]{spaces-divisors}
\externaldocument[spaces-over-fields-]{spaces-over-fields}
\externaldocument[spaces-topologies-]{spaces-topologies}
\externaldocument[spaces-descent-]{spaces-descent}
\externaldocument[spaces-perfect-]{spaces-perfect}
\externaldocument[spaces-more-morphisms-]{spaces-more-morphisms}
\externaldocument[spaces-flat-]{spaces-flat}
\externaldocument[spaces-groupoids-]{spaces-groupoids}
\externaldocument[spaces-more-groupoids-]{spaces-more-groupoids}
\externaldocument[bootstrap-]{bootstrap}
\externaldocument[spaces-pushouts-]{spaces-pushouts}
\externaldocument[spaces-chow-]{spaces-chow}
\externaldocument[groupoids-quotients-]{groupoids-quotients}
\externaldocument[spaces-more-cohomology-]{spaces-more-cohomology}
\externaldocument[spaces-simplicial-]{spaces-simplicial}
\externaldocument[spaces-duality-]{spaces-duality}
\externaldocument[formal-spaces-]{formal-spaces}
\externaldocument[restricted-]{restricted}
\externaldocument[spaces-resolve-]{spaces-resolve}
\externaldocument[formal-defos-]{formal-defos}
\externaldocument[defos-]{defos}
\externaldocument[cotangent-]{cotangent}
\externaldocument[examples-defos-]{examples-defos}
\externaldocument[algebraic-]{algebraic}
\externaldocument[examples-stacks-]{examples-stacks}
\externaldocument[stacks-sheaves-]{stacks-sheaves}
\externaldocument[criteria-]{criteria}
\externaldocument[artin-]{artin}
\externaldocument[quot-]{quot}
\externaldocument[stacks-properties-]{stacks-properties}
\externaldocument[stacks-morphisms-]{stacks-morphisms}
\externaldocument[stacks-limits-]{stacks-limits}
\externaldocument[stacks-cohomology-]{stacks-cohomology}
\externaldocument[stacks-perfect-]{stacks-perfect}
\externaldocument[stacks-introduction-]{stacks-introduction}
\externaldocument[stacks-more-morphisms-]{stacks-more-morphisms}
\externaldocument[stacks-geometry-]{stacks-geometry}
\externaldocument[moduli-]{moduli}
\externaldocument[moduli-curves-]{moduli-curves}
\externaldocument[examples-]{examples}
\externaldocument[exercises-]{exercises}
\externaldocument[guide-]{guide}
\externaldocument[desirables-]{desirables}
\externaldocument[coding-]{coding}
\externaldocument[obsolete-]{obsolete}
\externaldocument[fdl-]{fdl}
\externaldocument[index-]{index}

% Theorem environments.
%
\theoremstyle{plain}
\newtheorem{theorem}[subsection]{Theorem}
\newtheorem{proposition}[subsection]{Proposition}
\newtheorem{lemma}[subsection]{Lemma}

\theoremstyle{definition}
\newtheorem{definition}[subsection]{Definition}
\newtheorem{example}[subsection]{Example}
\newtheorem{exercise}[subsection]{Exercise}
\newtheorem{situation}[subsection]{Situation}

\theoremstyle{remark}
\newtheorem{remark}[subsection]{Remark}
\newtheorem{remarks}[subsection]{Remarks}

\numberwithin{equation}{subsection}

% Macros
%
\def\lim{\mathop{\mathrm{lim}}\nolimits}
\def\colim{\mathop{\mathrm{colim}}\nolimits}
\def\Spec{\mathop{\mathrm{Spec}}}
\def\Hom{\mathop{\mathrm{Hom}}\nolimits}
\def\Ext{\mathop{\mathrm{Ext}}\nolimits}
\def\SheafHom{\mathop{\mathcal{H}\!\mathit{om}}\nolimits}
\def\SheafExt{\mathop{\mathcal{E}\!\mathit{xt}}\nolimits}
\def\Sch{\mathit{Sch}}
\def\Mor{\mathop{\mathrm{Mor}}\nolimits}
\def\Ob{\mathop{\mathrm{Ob}}\nolimits}
\def\Sh{\mathop{\mathit{Sh}}\nolimits}
\def\NL{\mathop{N\!L}\nolimits}
\def\CH{\mathop{\mathrm{CH}}\nolimits}
\def\proetale{{pro\text{-}\acute{e}tale}}
\def\etale{{\acute{e}tale}}
\def\QCoh{\mathit{QCoh}}
\def\Ker{\mathop{\mathrm{Ker}}}
\def\Im{\mathop{\mathrm{Im}}}
\def\Coker{\mathop{\mathrm{Coker}}}
\def\Coim{\mathop{\mathrm{Coim}}}

% Boxtimes
%
\DeclareMathSymbol{\boxtimes}{\mathbin}{AMSa}{"02}

%
% Macros for moduli stacks/spaces
%
\def\QCohstack{\mathcal{QC}\!\mathit{oh}}
\def\Cohstack{\mathcal{C}\!\mathit{oh}}
\def\Spacesstack{\mathcal{S}\!\mathit{paces}}
\def\Quotfunctor{\mathrm{Quot}}
\def\Hilbfunctor{\mathrm{Hilb}}
\def\Curvesstack{\mathcal{C}\!\mathit{urves}}
\def\Polarizedstack{\mathcal{P}\!\mathit{olarized}}
\def\Complexesstack{\mathcal{C}\!\mathit{omplexes}}
% \Pic is the operator that assigns to X its picard group, usage \Pic(X)
% \Picardstack_{X/B} denotes the Picard stack of X over B
% \Picardfunctor_{X/B} denotes the Picard functor of X over B
\def\Pic{\mathop{\mathrm{Pic}}\nolimits}
\def\Picardstack{\mathcal{P}\!\mathit{ic}}
\def\Picardfunctor{\mathrm{Pic}}
\def\Deformationcategory{\mathcal{D}\!\mathit{ef}}
\usepackage{cleveref}
\crefname{originallemma}{original lemma}{original lemmas}
\Crefname{originallemma}{Original lemma}{Original lemmas}
\makeatletter
\newlabel{bootstrap-lemma-after-fppf-sep-lqf}{{046J}{source}{}{sourcepointer-1}{}}
\newlabel{bootstrap-lemma-after-fppf-sep-lqf@cref}{{[originallemma][0][]046J}{[1][0][]source}{}{}{}}
\makeatother
\crefname{originallemma}{original lemma}{original lemmas}
\Crefname{originallemma}{Original lemma}{Original lemmas}
\makeatletter
\newlabel{descent-lemma-reduced-local-smooth}{{034E}{source}{}{sourcepointer-2}{}}
\newlabel{descent-lemma-reduced-local-smooth@cref}{{[originallemma][0][]034E}{[1][0][]source}{}{}{}}
\makeatother
\crefname{originallemma}{original lemma}{original lemmas}
\Crefname{originallemma}{Original lemma}{Original lemmas}
\makeatletter
\newlabel{lemma-covering-by-thickenings}{{0AI9}{source}{}{sourcepointer-3}{}}
\newlabel{lemma-covering-by-thickenings@cref}{{[originallemma][0][]0AI9}{[1][0][]source}{}{}{}}
\makeatother
\crefname{originallemma}{original lemma}{original lemmas}
\Crefname{originallemma}{Original lemma}{Original lemmas}
\makeatletter
\newlabel{lemma-presentation-representable}{{0AIG}{source}{}{sourcepointer-4}{}}
\newlabel{lemma-presentation-representable@cref}{{[originallemma][0][]0AIG}{[1][0][]source}{}{}{}}
\makeatother
\crefname{originaldefinition}{original definition}{original definitions}
\Crefname{originaldefinition}{Original definition}{Original definitions}
\makeatletter
\newlabel{schemes-definition-reduced-induced-scheme}{{01J4}{source}{}{sourcepointer-5}{}}
\newlabel{schemes-definition-reduced-induced-scheme@cref}{{[originaldefinition][0][]01J4}{[1][0][]source}{}{}{}}
\makeatother
\crefname{originallemma}{original lemma}{original lemmas}
\Crefname{originallemma}{Original lemma}{Original lemmas}
\makeatletter
\newlabel{schemes-lemma-map-into-reduction}{{0356}{source}{}{sourcepointer-6}{}}
\newlabel{schemes-lemma-map-into-reduction@cref}{{[originallemma][0][]0356}{[1][0][]source}{}{}{}}
\makeatother
\crefname{originallemma}{original lemma}{original lemmas}
\Crefname{originallemma}{Original lemma}{Original lemmas}
\makeatletter
\newlabel{spaces-lemma-descent-representable-transformations-property}{{03KD}{source}{}{sourcepointer-7}{}}
\newlabel{spaces-lemma-descent-representable-transformations-property@cref}{{[originallemma][0][]03KD}{[1][0][]source}{}{}{}}
\makeatother
\crefname{originaldefinition}{original definition}{original definitions}
\Crefname{originaldefinition}{Original definition}{Original definitions}
\makeatletter
\newlabel{spaces-properties-definition-reduced-induced-space}{{047X}{source}{}{sourcepointer-8}{}}
\newlabel{spaces-properties-definition-reduced-induced-space@cref}{{[originaldefinition][0][]047X}{[1][0][]source}{}{}{}}
\makeatother
\crefname{originallemma}{original lemma}{original lemmas}
\Crefname{originallemma}{Original lemma}{Original lemmas}
\makeatletter
\newlabel{spaces-properties-lemma-map-into-reduction}{{03JJ}{source}{}{sourcepointer-9}{}}
\newlabel{spaces-properties-lemma-map-into-reduction@cref}{{[originallemma][0][]03JJ}{[1][0][]source}{}{}{}}
\makeatother
\crefname{originalsection}{original section}{original sections}
\Crefname{originalsection}{Original section}{Original sections}
\makeatletter
\newlabel{spaces-properties-section-types-properties}{{03E5}{source}{}{sourcepointer-10}{}}
\newlabel{spaces-properties-section-types-properties@cref}{{[originalsection][0][]03E5}{[1][0][]source}{}{}{}}
\makeatother
\crefname{originaltheorem}{original theorem}{original theorems}
\Crefname{originaltheorem}{Original theorem}{Original theorems}
\makeatletter
\newlabel{spaces-theorem-presentation}{{02WW}{source}{}{sourcepointer-11}{}}
\newlabel{spaces-theorem-presentation@cref}{{[originaltheorem][0][]02WW}{[1][0][]source}{}{}{}}
\makeatother
\crefname{originalremark}{original remark}{original remarks}
\Crefname{originalremark}{Original remark}{Original remarks}
\makeatletter
\newlabel{remark-set-theoretic}{{0AIS}{source}{}{sourcepointer-12}{}}
\newlabel{remark-set-theoretic@cref}{{[originalremark][0][]0AIS}{[1][0][]source}{}{}{}}
\makeatother
\title{A formal algebraic space has a reduced algebraic space with a representable thickening onto it}
\author{The Stacks Project Authors}
\date{stacks-project:formal-spaces.tex:local-snapshot-20261012; GNU FDL 1.2 or later}
\begin{document}
\maketitle
\noindent\textit{Editorial note.} All delivered mathematical text is the author's own, copied from the pinned Stacks Project chapter file formal-spaces.tex; no statement or proof step was rewritten, reordered or removed.\par\medskip
\noindent\textit{Editorial note.} The proof contains the author's sentence 'We omit the verification of the universal property (use ...)' followed by four references; it is reproduced verbatim and recorded as an obligation of the card. The universal property is nevertheless proved in full for the general case in the third part of the proof, so the omission concerns only the affine case.\par\medskip
\noindent\textit{Editorial note.} The statement refers to the notions of a representation of a formal algebraic space and of a thickening; both are the chapter's standing notions and their defining environments are not reproduced, the references being carried as navigation stubs.\par\medskip
\noindent\textit{Editorial note.} The proof displays a cartesian square and a chain of identifications of fibre products; both lie inside the counted span and are reproduced verbatim.\par\medskip
\noindent\textit{Editorial note.} The counted spans carry 13 navigation stubs; no bibliography citation occurs in the spans.\par\medskip
\noindent\textit{Editorial note.} The preamble span points at the chapter's own preamble.tex shipped with the frozen source, used verbatim as staged.\par\medskip
\noindent\textit{Allowed standard prerequisites (editorial).}\begin{itemize}
\item Affine formal algebraic spaces as filtered colimits of affine schemes, and the thickenings between them. Same chapter, Definition definition-affine-formal-algebraic-space, invoked at proof lines 2841-2842. Supplies the system whose stages define the reduction in the affine case.
\item Being a thickening is stable under composition, and the reduction of a scheme with its universal property. Same chapter, Lemma lemma-covering-by-thickenings at proof lines 2846-2847, and Schemes, Definition schemes-definition-reduced-induced-scheme and Lemma schemes-lemma-map-into-reduction, listed at proof lines 2850-2852. Gives representability and the thickening property in the affine case and the model for the universal property.
\item The definitions of reduced algebraic spaces and of the reduction of an algebraic space. Properties of Spaces, Definition spaces-properties-definition-reduced-induced-space and Lemma spaces-properties-lemma-map-into-reduction, listed at proof lines 2852-2854, and Section spaces-properties-section-types-properties at proof lines 2895-2897. Supplies what a reduced algebraic space is and the reduction of an algebraic space with its functorial property.
\item Etale morphisms of schemes have reduced sources over reduced targets. Descent, Lemma descent-lemma-reduced-local-smooth, invoked at proof lines 2862-2864. Makes the fibre products appearing in the glueing construction reduced, so that they factor through the reductions.
\item Etale equivalence relations present algebraic spaces, and representability descends along a surjective etale morphism. Spaces, Theorem spaces-theorem-presentation at proof lines 2878-2879, Bootstrap, Lemma bootstrap-lemma-after-fppf-sep-lqf at proof lines 2884-2886, and Spaces, Lemma spaces-lemma-descent-representable-transformations-property at proof lines 2887-2888. Produces the quotient algebraic space and transfers representability, and then the thickening property, from the affine pieces to the whole.
\end{itemize}\medskip
\ifdefined\hypertarget\hypertarget{sourcepointer-1}{}\fi
\noindent\textit{Source reference.} Original lemma 046J: sources/stacks-chow/tags.txt, source lines 5474--5474. Pointer outside the counted span, invoked once for context; the load-bearing argument of the counted statement is the author's proof reproduced in full in the proof span.\par\medskip
\ifdefined\hypertarget\hypertarget{sourcepointer-2}{}\fi
\noindent\textit{Source reference.} Original lemma 034E: sources/stacks-chow/tags.txt, source lines 4089--4089. Pointer outside the counted span, invoked once for context; the load-bearing argument of the counted statement is the author's proof reproduced in full in the proof span.\par\medskip
\ifdefined\hypertarget\hypertarget{sourcepointer-3}{}\fi
\noindent\textit{Source reference.} Original lemma 0AI9: sources/stacks-formal-spaces/pinned-original-source.tex, source lines 2057--2057. Pointer outside the counted span, invoked once for context; the load-bearing argument of the counted statement is the author's proof reproduced in full in the proof span.\par\medskip
\ifdefined\hypertarget\hypertarget{sourcepointer-4}{}\fi
\noindent\textit{Source reference.} Original lemma 0AIG: sources/stacks-formal-spaces/pinned-original-source.tex, source lines 2381--2381. Pointer outside the counted span, invoked once for context; the load-bearing argument of the counted statement is the author's proof reproduced in full in the proof span.\par\medskip
\ifdefined\hypertarget\hypertarget{sourcepointer-5}{}\fi
\noindent\textit{Source reference.} Original definition 01J4: sources/stacks-chow/tags.txt, source lines 2001--2001. Pointer outside the counted span, invoked once for context; the load-bearing argument of the counted statement is the author's proof reproduced in full in the proof span.\par\medskip
\ifdefined\hypertarget\hypertarget{sourcepointer-6}{}\fi
\noindent\textit{Source reference.} Original lemma 0356: sources/stacks-chow/tags.txt, source lines 4117--4117. Pointer outside the counted span, invoked once for context; the load-bearing argument of the counted statement is the author's proof reproduced in full in the proof span.\par\medskip
\ifdefined\hypertarget\hypertarget{sourcepointer-7}{}\fi
\noindent\textit{Source reference.} Original lemma 03KD: sources/stacks-chow/tags.txt, source lines 4673--4673. Pointer outside the counted span, invoked once for context; the load-bearing argument of the counted statement is the author's proof reproduced in full in the proof span.\par\medskip
\ifdefined\hypertarget\hypertarget{sourcepointer-8}{}\fi
\noindent\textit{Source reference.} Original definition 047X: sources/stacks-chow/tags.txt, source lines 5524--5524. Pointer outside the counted span, invoked once for context; the load-bearing argument of the counted statement is the author's proof reproduced in full in the proof span.\par\medskip
\ifdefined\hypertarget\hypertarget{sourcepointer-9}{}\fi
\noindent\textit{Source reference.} Original lemma 03JJ: sources/stacks-chow/tags.txt, source lines 4641--4641. Pointer outside the counted span, invoked once for context; the load-bearing argument of the counted statement is the author's proof reproduced in full in the proof span.\par\medskip
\ifdefined\hypertarget\hypertarget{sourcepointer-10}{}\fi
\noindent\textit{Source reference.} Original section 03E5: sources/stacks-chow/tags.txt, source lines 4444--4444. Pointer outside the counted span, invoked once for context; the load-bearing argument of the counted statement is the author's proof reproduced in full in the proof span.\par\medskip
\ifdefined\hypertarget\hypertarget{sourcepointer-11}{}\fi
\noindent\textit{Source reference.} Original theorem 02WW: sources/stacks-chow/tags.txt, source lines 3819--3819. Pointer outside the counted span, invoked once for context; the load-bearing argument of the counted statement is the author's proof reproduced in full in the proof span.\par\medskip
\ifdefined\hypertarget\hypertarget{sourcepointer-12}{}\fi
\noindent\textit{Source reference.} Original remark 0AIS: sources/stacks-formal-spaces/pinned-original-source.tex, source lines 2769--2769. Pointer outside the counted span, invoked once for context; the load-bearing argument of the counted statement is the author's proof reproduced in full in the proof span.\par\medskip
\begin{definition}
\label{definition-affine-formal-algebraic-space}
Let $S$ be a scheme. We say a sheaf $X$ on $(\Sch/S)_{fppf}$ is an
{\it affine formal algebraic space} if there exist
\begin{enumerate}
\item a directed set $\Lambda$,
\item a system $(X_\lambda, f_{\lambda \mu})$ over $\Lambda$
in $(\Sch/S)_{fppf}$ where
\begin{enumerate}
\item each $X_\lambda$ is affine,
\item each $f_{\lambda \mu} : X_\lambda \to X_\mu$ is a thickening,
\end{enumerate}
\end{enumerate}
such that
$$
X \cong \colim_{\lambda \in \Lambda} X_\lambda
$$
as fppf sheaves and $X$ satisfies a set theoretic condition
(see Remark \ref{remark-set-theoretic}). A
{\it morphism of affine formal algebraic spaces}
over $S$ is a map of sheaves.
\end{definition}

\begin{definition}
\label{definition-formal-algebraic-space}
Let $S$ be a scheme. We say a sheaf $X$ on $(\Sch/S)_{fppf}$ is a
{\it formal algebraic space} if there exist a family of maps
$\{X_i \to X\}_{i \in I}$ of sheaves such that
\begin{enumerate}
\item $X_i$ is an affine formal algebraic space,
\item $X_i \to X$ is representable by algebraic spaces and \'etale,
\item $\coprod X_i \to X$ is surjective as a map of sheaves
\end{enumerate}
and $X$ satisfies a set theoretic condition
(see Remark \ref{remark-set-theoretic}). A
{\it morphism of formal algebraic spaces}
over $S$ is a map of sheaves.
\end{definition}

\begin{lemma}
\label{lemma-reduction-formal-algebraic-space}
Let $S$ be a scheme. Let $X$ be a formal algebraic space over $S$.
There exists a reduced algebraic space $X_{red}$ and a representable
morphism $X_{red} \to X$ which is a thickening. A morphism $U \to X$
with $U$ a reduced algebraic space factors uniquely through $X_{red}$.
\end{lemma}

\begin{proof}
First assume that $X$ is an affine formal algebraic space.
Say $X = \colim X_\lambda$ as in
Definition \ref{definition-affine-formal-algebraic-space}.
Since the transition morphisms are thickenings, the affine
schemes $X_\lambda$ all have isomorphic reductions $X_{red}$.
The morphism $X_{red} \to X$ is representable and a thickening
by Lemma \ref{lemma-covering-by-thickenings} and the fact that
compositions of thickenings are thickenings. We omit the
verification of the universal
property (use Schemes, Definition
\ref{schemes-definition-reduced-induced-scheme},
Schemes, Lemma \ref{schemes-lemma-map-into-reduction},
Properties of Spaces, Definition
\ref{spaces-properties-definition-reduced-induced-space}, and
Properties of Spaces, Lemma \ref{spaces-properties-lemma-map-into-reduction}).

\medskip\noindent
Let $X$ and $\{X_i \to X\}_{i \in I}$ be as in
Definition \ref{definition-formal-algebraic-space}.
For each $i$ let $X_{i, red} \to X_i$ be the reduction as
constructed above. For $i, j \in I$ the projection
$X_{i, red} \times_X X_j \to X_{i, red}$ is an \'etale (by assumption)
morphism of schemes (by Lemma \ref{lemma-presentation-representable}).
Hence $X_{i, red} \times_X X_j$ is reduced (see
Descent, Lemma \ref{descent-lemma-reduced-local-smooth}).
Thus the projection $X_{i, red} \times_X X_j \to X_j$ factors
through $X_{j, red}$ by the universal property. We conclude that
$$
R_{ij} = X_{i, red} \times_X X_j = X_{i, red} \times_X X_{j, red} =
X_i \times_X X_{j, red}
$$
because the morphisms $X_{i, red} \to X_i$ are injections of sheaves.
Set $U = \coprod X_{i, red}$, set
$R = \coprod R_{ij}$, and denote $s, t : R \to U$ the two
projections. As a sheaf $R = U \times_X U$ and $s$ and $t$
are \'etale. Then $(t, s) : R \to U$ defines an \'etale equivalence
relation by our observations above. Thus $X_{red} = U/R$ is an
algebraic space by Spaces, Theorem \ref{spaces-theorem-presentation}.
By construction the diagram
$$
\xymatrix{
\coprod X_{i, red} \ar[r] \ar[d] & \coprod X_i \ar[d] \\
X_{red} \ar[r] & X
}
$$
is cartesian. Since the right vertical arrow is \'etale surjective
and the top horizontal arrow is representable and a thickening
we conclude that $X_{red} \to X$ is representable by
Bootstrap, Lemma \ref{bootstrap-lemma-after-fppf-sep-lqf}
(to verify the assumptions of the lemma use that a surjective
\'etale morphism is surjective, flat, and locally of finite
presentation and use that thickenings are separated and locally quasi-finite).
Then we can use Spaces, Lemma
\ref{spaces-lemma-descent-representable-transformations-property}
to conclude that $X_{red} \to X$ is a thickening
(use that being a thickening is equivalent to being
a surjective closed immersion).

\medskip\noindent
Finally, suppose that $U \to X$ is a morphism with
$U$ a reduced algebraic space over $S$. Then each $X_i \times_X U$
is \'etale over $U$ and therefore reduced (by our definition of
reduced algebraic spaces in Properties of Spaces, Section
\ref{spaces-properties-section-types-properties}).
Then $X_i \times_X U \to X_i$ factors through $X_{i, red}$.
Hence $U \to X$ factors through $X_{red}$ because
$\{X_i \times_X U \to U\}$ is an \'etale covering.
\end{proof}

\end{document}
